\documentclass[11pt]{article}

\usepackage[a4paper,margin=1in]{geometry}
\usepackage[T1]{fontenc}
\usepackage[utf8]{inputenc}
\usepackage{amsmath,amssymb,amsfonts}
\usepackage{booktabs}
\usepackage{longtable}
\usepackage{array}
\usepackage{enumitem}
\usepackage{hyperref}

\hypersetup{
  colorlinks=true,
  linkcolor=blue,
  urlcolor=blue,
  citecolor=blue
}

\title{Designing a Reward / Cost Function for Multi-Objective Unitary Synthesis in \texttt{gqe-torch}}
\author{}
\date{\today}

\begin{document}
\maketitle
\tableofcontents

\section{Purpose}

This document is a \emph{design proposal} for a reward / cost function suitable for
the current \texttt{gqe-torch} training setup.
The goal is to produce a scalar objective that:

\begin{itemize}[leftmargin=*]
\item strongly prioritizes process fidelity early in training;
\item increasingly favors low depth and low CNOT count once fidelity is already good;
\item preserves an exponential / logarithmic coupling between fidelity and structure;
\item avoids numerical pathologies at $F=0$ and $F=1$;
\item remains compatible with replay-based GRPO training.
\end{itemize}

The core problem is multi-objective:

\begin{itemize}[leftmargin=*]
\item maximize fidelity $F$;
\item minimize circuit depth $D$;
\item minimize CNOT count $C$.
\end{itemize}

In the current codebase, total gate count is fixed by construction, so it is not a useful
training objective yet.

\section{Project Constraints}

Any reward / cost used here should respect the following properties of the current system:

\begin{enumerate}[leftmargin=*]
\item The generator emits \emph{fixed-length} token sequences.
\item The replay buffer stores scalar costs, and those costs are reused across later updates.
\item The training loop is not using a stepwise environment reward; it is scoring complete circuits.
\item Reference quantities that move too quickly can make replay targets inconsistent.
\item Small differences near $F \approx 1$ matter much more than similar absolute differences near $F \approx 0.7$.
\end{enumerate}

These constraints mean the training objective should be as close to \emph{stationary} as possible.
For that reason, a pure ``reward = change from previous epoch'' formulation is not ideal for
replay training, even if it is intuitive for human interpretation.

\section{Why the Naive Exponential Proposal Is Risky}

The intuitive form
\[
r
=
\left(
\alpha \frac{D^* - 1}{D_{\mathrm{current}}}
+
\beta \frac{C^* - 1}{C_{\mathrm{current}}}
\right)^{\frac{1}{1-F}}
\]
captures an important idea:
\emph{structure improvements should matter more when fidelity is already high.}

That idea is good, but the raw formula has several problems.

\subsection{Singularity at \texorpdfstring{$F=1$}{F=1}}

The exponent $\frac{1}{1-F}$ diverges as $F \to 1$.
This creates an unbounded numerical instability exactly in the regime where the training
signal needs to be most reliable.

\subsection{Exploding or Vanishing Reward}

Let
\[
B
=
\alpha \frac{D^* - 1}{D_{\mathrm{current}}}
+
\beta \frac{C^* - 1}{C_{\mathrm{current}}}.
\]

If $B > 1$, then $B^{1/(1-F)}$ can explode.
If $0 < B < 1$, then $B^{1/(1-F)}$ can collapse toward zero extremely fast.

So the fidelity coupling is too aggressive and not well controlled.

\subsection{Bad Behavior When \texorpdfstring{$D^*=1$}{D*=1} or \texorpdfstring{$C^*=1$}{C*=1}}

The terms $D^*-1$ and $C^*-1$ can become zero.
That can collapse the whole base even when the current circuit is good.

\subsection{Poisoned References}

If $D^*$ and $C^*$ are taken as unconstrained ``best so far'' values from \emph{all} circuits,
then a low-fidelity but extremely compact circuit can make the structure target unrealistically hard.
This is exactly the kind of behavior that can push the model to over-optimize structure while sacrificing fidelity.

\subsection{Conclusion}

The right design choice is therefore:

\begin{itemize}[leftmargin=*]
\item keep the \emph{exponential / logarithmic coupling};
\item replace the singular exponent by a \emph{bounded, fidelity-dependent exponent};
\item define structure references carefully;
\item separate a \emph{stable training cost} from an optional \emph{signed diagnostic reward}.
\end{itemize}

\section{Proposed Design}

\subsection{Notation}

Let:

\begin{itemize}[leftmargin=*]
\item $F \in [0,1]$ be process fidelity;
\item $D \in \{1,2,\dots\}$ be estimated circuit depth;
\item $C \in \{0,1,2,\dots\}$ be CNOT count;
\item $\alpha,\beta \ge 0$ with $\alpha+\beta=1$ be the depth / CNOT blending weights.
\end{itemize}

We also define:

\begin{itemize}[leftmargin=*]
\item $\varepsilon_F > 0$ as a small numerical constant for fidelity shaping;
\item $\varepsilon_B > 0$ as a small numerical constant for the structure term;
\item $\delta_D = \delta_C = 1$ as offsets preventing division by zero.
\end{itemize}

\subsection{Reference Depth and CNOT Counts}

We should \emph{not} define $D^*$ and $C^*$ as the best values ever seen regardless of fidelity.
Instead, they should come from a fidelity-qualified set.

Choose a reference fidelity threshold $\tau_{\mathrm{ref}}$ and define
\[
\mathcal{A}_{\tau_{\mathrm{ref}}}
=
\left\{
(F_i,D_i,C_i)
\;:\;
F_i \ge \tau_{\mathrm{ref}}
\right\}.
\]

Then define
\[
D^*
=
\begin{cases}
\min\{D_i : (F_i,D_i,C_i)\in \mathcal{A}_{\tau_{\mathrm{ref}}}\}, & \mathcal{A}_{\tau_{\mathrm{ref}}}\neq\varnothing,\\[6pt]
D_{\mathrm{fallback}}, & \mathcal{A}_{\tau_{\mathrm{ref}}}=\varnothing,
\end{cases}
\]
\[
C^*
=
\begin{cases}
\min\{C_i : (F_i,D_i,C_i)\in \mathcal{A}_{\tau_{\mathrm{ref}}}\}, & \mathcal{A}_{\tau_{\mathrm{ref}}}\neq\varnothing,\\[6pt]
C_{\mathrm{fallback}}, & \mathcal{A}_{\tau_{\mathrm{ref}}}=\varnothing.
\end{cases}
\]

Good fallback choices are:

\begin{itemize}[leftmargin=*]
\item the depth and CNOT count of the best-fidelity circuit seen so far;
\item a slowly updated EMA reference;
\item or the best reference values obtained during warmup.
\end{itemize}

\subsection{Fidelity Shaping}

We want fidelity to dominate early, and we want high-fidelity circuits to be well separated.
Define
\[
\phi(F)
=
\operatorname{clip}\!\left(
-\log(1-F+\varepsilon_F),
\,
0,
\,
\phi_{\max}
\right).
\]

Properties:

\begin{itemize}[leftmargin=*]
\item $\phi(F)$ is small when $F$ is poor;
\item $\phi(F)$ grows rapidly as $F \to 1$;
\item clipping prevents blow-up.
\end{itemize}

\subsection{Structure Ratios}

Define structure-efficiency ratios
\[
R_D(D)
=
\frac{D^* + \delta_D}{D + \delta_D},
\qquad
R_C(C)
=
\frac{C^* + \delta_C}{C + \delta_C}.
\]

These satisfy:

\begin{itemize}[leftmargin=*]
\item $R_D = 1$ when $D=D^*$;
\item $R_D > 1$ when $D < D^*$;
\item $R_D < 1$ when $D > D^*$;
\item and similarly for $R_C$.
\end{itemize}

Blend them into one structure score:
\[
B(D,C)
=
\alpha R_D(D)
+
\beta R_C(C).
\]

Thus:

\begin{itemize}[leftmargin=*]
\item $B=1$ means the circuit matches the reference structure;
\item $B>1$ means it is structurally better than the reference;
\item $B<1$ means it is structurally worse.
\end{itemize}

\subsection{Bounded Fidelity-Dependent Exponent}

To keep the ``structure matters more as fidelity improves'' idea, define a bounded exponent:
\[
\kappa(F)
=
\kappa_{\min}
+
\left(\kappa_{\max}-\kappa_{\min}\right)
\left[
\operatorname{clip}\!\left(
\frac{F-F_{\mathrm{start}}}{F_{\mathrm{full}}-F_{\mathrm{start}}},
\,
0,
\,
1
\right)
\right]^p.
\]

Interpretation:

\begin{itemize}[leftmargin=*]
\item for $F \le F_{\mathrm{start}}$, $\kappa(F)\approx \kappa_{\min}$;
\item for $F \ge F_{\mathrm{full}}$, $\kappa(F)\approx \kappa_{\max}$;
\item $p>1$ makes the ramp delayed and sharper.
\end{itemize}

This preserves the exponential idea without using the unstable factor $\frac{1}{1-F}$.

\subsection{Absolute Training Score}

Define the absolute circuit score
\[
S(F,D,C)
=
\left(\phi(F)+\varepsilon_B\right)\,
\left(B(D,C)+\varepsilon_B\right)^{\kappa(F)}.
\]

Equivalently, in log-space,
\[
s(F,D,C)
=
\log\!\left(\phi(F)+\varepsilon_B\right)
+
\kappa(F)\log\!\left(B(D,C)+\varepsilon_B\right).
\]

The log form is usually the numerically safer representation.

\subsection{Training Cost}

For replay-based GRPO training, use a cost rather than a signed improvement reward:
\[
C_{\mathrm{train}}(F,D,C)
=
-\,s(F,D,C).
\]

Lower cost is better.

This is the key recommendation for the actual training loop:
\begin{itemize}[leftmargin=*]
\item use a \emph{stationary or slowly moving absolute cost} for replay;
\item do not use raw epoch-to-epoch reward differences as the replay target.
\end{itemize}

\subsection{Optional Signed Diagnostic Reward}

For analysis, logging, or acceptance rules, define a signed reward relative to a reference circuit:
\[
r_{\mathrm{diag}}(F,D,C)
=
\bigl[\phi(F)-\phi(F_{\mathrm{ref}})\bigr]
+
\mathbf{1}\!\left[F \ge F_{\mathrm{ref}}-\Delta_F\right]
\eta\,\kappa(F)\log\!\left(B(D,C)+\varepsilon_B\right).
\]

Here:

\begin{itemize}[leftmargin=*]
\item $(F_{\mathrm{ref}},D_{\mathrm{ref}},C_{\mathrm{ref}})$ is a chosen anchor circuit;
\item $\Delta_F \ge 0$ is a fidelity-drop tolerance;
\item $\eta > 0$ scales the structure contribution in the diagnostic reward.
\end{itemize}

This diagnostic reward has the behavior many users expect intuitively:

\begin{itemize}[leftmargin=*]
\item if fidelity improves and structure is similar, reward is positive;
\item if fidelity is unchanged and structure improves, reward is positive;
\item if fidelity decreases by more than $\Delta_F$, the structure term is suppressed and reward becomes negative.
\end{itemize}

\section{Recommended Interpretation}

The proposed design therefore uses \emph{two related but distinct quantities}:

\begin{enumerate}[leftmargin=*]
\item \textbf{Absolute training cost} $C_{\mathrm{train}}$ for replay-based optimization.
\item \textbf{Signed diagnostic reward} $r_{\mathrm{diag}}$ for interpretation and policy analysis.
\end{enumerate}

This separation is important.
The signed reward is easy to reason about, but the absolute cost is more suitable for stable replay training.

\section{Edge-Case Analysis}

\begin{longtable}{>{\raggedright\arraybackslash}p{0.23\linewidth} >{\raggedright\arraybackslash}p{0.71\linewidth}}
\toprule
Case & Behavior of the proposed design \\
\midrule
\endfirsthead
\toprule
Case & Behavior of the proposed design \\
\midrule
\endhead
\midrule
\multicolumn{2}{r}{Continued on next page}
\endfoot
\bottomrule
\endlastfoot

$F=0$ &
$\phi(F)\approx 0$, so $S(F,D,C)$ is near zero and $C_{\mathrm{train}}$ is large.
Even very good structure cannot rescue a useless circuit. This is desirable. \\[4pt]

$0<F<F_{\mathrm{start}}$ &
The exponent $\kappa(F)$ stays close to $\kappa_{\min}$, so structure matters only weakly.
Training is fidelity-first. \\[4pt]

$F=F_{\mathrm{start}}$ &
This is where structure pressure begins to ramp.
The reward does not jump discontinuously; it changes smoothly. \\[4pt]

$F_{\mathrm{start}}<F<F_{\mathrm{full}}$ &
This is the trade-off region.
Structure gains increasingly matter, but fidelity still has strong influence through $\phi(F)$. \\[4pt]

$F\to 1$ &
$\phi(F)$ becomes large but clipped, and $\kappa(F)\to \kappa_{\max}$.
High-fidelity circuits are then separated strongly by depth and CNOT count.
There is no singularity. \\[4pt]

$F=1$ exactly &
The design remains finite because $\phi(F)$ uses $\varepsilon_F$ and clipping.
Unlike $1/(1-F)$, nothing diverges. \\[4pt]

$D=D^*, C=C^*$ &
Then $R_D=R_C=1$ and $B=1$.
The structure term is neutral.
The score reduces to a fidelity-only shaped score. \\[4pt]

$D<D^*$, $C=C^*$ &
Then $R_D>1$, so $B>1$ and the score increases.
If the circuit is kept, the reference depth should eventually be updated. \\[4pt]

$D=D^*$, $C<C^*$ &
Then $R_C>1$, so $B>1$ and the score increases. \\[4pt]

$D<D^*$ and $C<C^*$ &
Both structure ratios exceed one, so the structure multiplier is strongly favorable. \\[4pt]

$D>D^*$ and $C>C^*$ &
Both structure ratios are below one, so the structure multiplier penalizes the circuit. \\[4pt]

Fidelity decreases slightly, structure improves slightly &
For the \emph{training cost}, the answer depends on the full trade-off.
For the \emph{diagnostic reward}, if the fidelity drop exceeds $\Delta_F$, the reward is forced negative. If the drop is within tolerance, the structure gain may still partially compensate. \\[4pt]

Fidelity improves, structure worsens slightly &
Early in training this is usually still favorable because fidelity dominates.
Late in training the sign can flip if the structure regression is too large. \\[4pt]

No fidelity-qualified reference exists yet &
Use fallback references from warmup, best-fidelity-so-far, or a slowly moving EMA.
This prevents the structure term from being defined by low-fidelity compact circuits. \\[4pt]

$D^*=1$ or $C^*=0$ &
The $+1$ offsets keep $R_D$ and $R_C$ well defined.
The reward does not collapse to zero. \\[4pt]

$C=0$ &
Safe because of the offset $\delta_C=1$.
If zero-CNOT circuits are physically plausible for the target, they are rewarded appropriately. \\[4pt]

Current circuit equals the reference in fidelity but has better structure &
The absolute training score improves through $B>1$.
The diagnostic reward is positive because the structure term is positive while the fidelity difference term is zero. \\[4pt]

Current circuit has lower fidelity than the reference but much better structure &
This is the dangerous case.
The design handles it by keeping $\kappa(F)$ small before high-fidelity regimes and by optionally suppressing the structure term in $r_{\mathrm{diag}}$ when fidelity drops by more than $\Delta_F$. \\[4pt]

\end{longtable}

\section{Why This Design Fits the Project}

This proposal matches the project better than a raw exponential in $\frac{1}{1-F}$ for several reasons:

\begin{itemize}[leftmargin=*]
\item it preserves the key intuition that structure should matter more at high fidelity;
\item it uses bounded, controlled nonlinearities instead of singular ones;
\item it is compatible with replay-based GRPO training through an absolute cost;
\item it naturally supports Pareto-style archive logic and high-fidelity structure comparisons;
\item it avoids rewarding low-fidelity but compact circuits too strongly.
\end{itemize}

\section{Recommended Hyperparameter Ranges}

Reasonable starting ranges are:

\begin{itemize}[leftmargin=*]
\item $\varepsilon_F = 10^{-8}$;
\item $\phi_{\max} \in [10,14]$;
\item $\kappa_{\min} \in [0.05,0.20]$;
\item $\kappa_{\max} \in [1.5,3.0]$;
\item $F_{\mathrm{start}} \in [0.90,0.97]$;
\item $F_{\mathrm{full}} \in [0.99,0.999]$;
\item $p \in [4,8]$;
\item $\Delta_F \in [10^{-4},10^{-3}]$;
\item $\alpha,\beta$ chosen so that the more important hardware bottleneck receives larger weight.
\end{itemize}

For three-qubit synthesis in a native \texttt{rz/sx/cx} basis, it is usually reasonable to set
$\beta > \alpha$ because CNOT count often dominates hardware cost.

\section{Final Recommendation}

For this repository, the most defensible design is:

\begin{enumerate}[leftmargin=*]
\item use a fidelity-shaped term $\phi(F)$;
\item use structure ratios relative to \emph{fidelity-qualified} best-so-far references;
\item use a bounded, delayed fidelity-dependent exponent $\kappa(F)$;
\item train with the absolute cost
\[
C_{\mathrm{train}}(F,D,C) = -\log(\phi(F)+\varepsilon_B) - \kappa(F)\log(B(D,C)+\varepsilon_B);
\]
\item optionally monitor the signed diagnostic reward
\[
r_{\mathrm{diag}}(F,D,C)
=
\bigl[\phi(F)-\phi(F_{\mathrm{ref}})\bigr]
+
\mathbf{1}[F\ge F_{\mathrm{ref}}-\Delta_F]\eta\kappa(F)\log(B(D,C)+\varepsilon_B).
\]
\end{enumerate}

This keeps the desirable exponential / logarithmic structure, but removes the instability
of the naive $1/(1-F)$ exponent.

\end{document}
